\section{Sujet n\textsuperscript{o}\,21}
\noindent\begin{minipage}{0.5\linewidth}\footnotesize Office du Baccalauréat du Cameroun\end{minipage}%
\hfill\begin{minipage}{0.45\linewidth}\footnotesize\raggedleft Examen : \textbf{Baccalauréat} · Série : \textbf{C/E}\\
Épreuve : \textbf{Mathématiques} · Durée : \textbf{4\,h} · Coef. : \textbf{7}\end{minipage}
\par\smallskip{\color{onbo}\rule{\linewidth}{1pt}}

\subsection*{Partie A — Évaluation des ressources \hfill (15 points)}

\noindent\textbf{Exercice 1.} \hfill\textit{(3 points)}\\
Un studio de musique de Douala dispose d'un catalogue de $10$ titres distincts pour
composer ses playlists.
\begin{enumerate}[label=\arabic*)]
\item Combien de playlists \emph{ordonnées} de $4$ titres distincts peut-on former ? \hfill\textit{(1 pt)}
\item On veut une playlist ordonnée de $4$ titres distincts \emph{commençant} par un titre
imposé (un jingle fixé du catalogue). Combien y en a-t-il ? \hfill\textit{(1 pt)}
\item Combien de sélections \emph{non ordonnées} de $4$ titres distincts peut-on former ? \hfill\textit{(1 pt)}
\end{enumerate}

\smallskip
\noindent\textbf{Exercice 2.} \hfill\textit{(3 points)}\\
Le studio relie ses quatre salles $A$, $B$, $C$, $D$ par des câbles audio orientés. La
matrice d'adjacence de ce réseau, dans l'ordre $(A,B,C,D)$, est
$\displaystyle M=\begin{pmatrix}0&1&1&0\\0&0&1&1\\0&0&0&1\\1&0&0&0\end{pmatrix}$
(le terme $(i,j)$ vaut $1$ s'il existe un câble direct de la salle $i$ vers la salle $j$).
\begin{enumerate}[label=\arabic*)]
\item Calculer $M^2$. \hfill\textit{(1,5 pt)}
\item En déduire le nombre de chemins de longueur $2$ allant de $A$ vers $D$, et les citer. \hfill\textit{(1 pt)}
\item Donner le nombre total de chemins de longueur $2$ dans le réseau. \hfill\textit{(0,5 pt)}
\end{enumerate}

\smallskip
\noindent\textbf{Exercice 3.} \hfill\textit{(4 points)}\\
Soit $f$ la fonction définie sur $\R$ par $f(x)=x\,\mathrm{e}^{-x}$, de courbe $(\mathcal{C})$
dans un repère orthonormé d'unité \SI{1}{\centi\metre}.
\begin{enumerate}[label=\arabic*)]
\item Déterminer les limites de $f$ en $-\infty$ et $+\infty$, et dresser le tableau de
variations de $f$. \hfill\textit{(1,5 pt)}
\item À l'aide d'une intégration par parties, calculer l'aire $\mathcal{A}$, en \si{\centi\metre\squared},
du domaine compris entre $(\mathcal{C})$, l'axe des abscisses et les droites $x=0$ et $x=1$. \hfill\textit{(1,5 pt)}
\item On considère $g(x)=\mathrm{e}^{-x}$. Calculer le volume $V$, en unités de volume,
du solide engendré par la rotation autour de l'axe des abscisses de la courbe de $g$ sur
$[0\,;1]$, c'est-à-dire $V=\pi\displaystyle\int_0^1 g(x)^2\dd x$. \hfill\textit{(1 pt)}
\end{enumerate}

\smallskip
\noindent\textbf{Exercice 4.} \hfill\textit{(5 points)}\\
Dans le plan muni d'un repère orthonormé, on considère les points $A(1\,;\,2)$ et $B(5\,;\,4)$.
\begin{enumerate}[label=\arabic*)]
\item Calculer $\overrightarrow{AB}\cdot\overrightarrow{AB}$ puis la distance $AB$. \hfill\textit{(1 pt)}
\item Déterminer et caractériser l'ensemble $(\Gamma_1)$ des points $M$ du plan tels que
$\overrightarrow{MA}\cdot\overrightarrow{MB}=0$. \hfill\textit{(1,5 pt)}
\item On note $I$ le milieu de $[AB]$. Déterminer l'ensemble $(\Gamma_2)$ des points $M$
tels que $MA^2+MB^2=40$ (on pourra utiliser $MA^2+MB^2=2MI^2+\tfrac12 AB^2$). \hfill\textit{(2,5 pt)}
\end{enumerate}

\subsection*{Partie B — Évaluation des compétences \hfill (5 points)}

\noindent\textit{Monsieur Nkoulou dirige le studio de production audiovisuelle de Douala
et te sollicite, en tant qu'élève de Terminale~C, pour trois décisions. Pour une séance,
le coût total de production, en milliers de francs CFA, en fonction du nombre $x$ de
titres enregistrés ($x\ge1$) est $C(x)=x^2-12x+80$, et chaque titre est vendu
\SI{10}{} milliers de francs. Par ailleurs, le nombre d'abonnés à sa chaîne, de $2000$
le premier mois, augmente de \SI{5}{\percent} chaque mois. Enfin, parmi les fichiers
livrés, \SI{8}{\percent} comportent une erreur d'encodage ; un client en télécharge $3$
au hasard.}

\begin{enumerate}[label=\textbf{Tâche \arabic*.},leftmargin=2.4em]
\item Déterminer le nombre de titres $x$ à enregistrer pour que le \emph{bénéfice} soit
maximal, et calculer ce bénéfice maximal. \hfill\textit{(1,5 pt)}
\item Déterminer le mois à partir duquel le nombre d'abonnés dépassera $3000$. \hfill\textit{(1,5 pt)}
\item Déterminer la probabilité qu'\emph{au moins un} des $3$ fichiers téléchargés comporte
une erreur d'encodage. \hfill\textit{(1,5 pt)}
\end{enumerate}
\par\smallskip{\footnotesize\itshape Présentation générale : 0,5 point.}

% ════════════════════════════════════════════════════
\corrige{du sujet 21}

\noindent\textbf{Partie A — Exercice 1.}
1) Playlists ordonnées de $4$ titres distincts parmi $10$ : arrangements
$A_{10}^4=10\times9\times8\times7=5040$.
2) Le premier titre est imposé ; il reste à ordonner $3$ titres distincts parmi les $9$
restants : $A_9^3=9\times8\times7=504$.
3) Sélections non ordonnées : $\binom{10}{4}=\frac{10\times9\times8\times7}{24}=210$.

\noindent\textbf{Exercice 2.}
1) $M^2=\begin{pmatrix}0&1&1&0\\0&0&1&1\\0&0&0&1\\1&0&0&0\end{pmatrix}^2
=\begin{pmatrix}0&0&1&2\\1&0&0&1\\1&0&0&0\\0&1&1&0\end{pmatrix}$.
\emph{Détail de la 1\textsuperscript{re} ligne} (ligne de $A=(0,1,1,0)$) : colonne $A$ :
$0$ ; colonne $B$ : $0\cdot1+1\cdot0+1\cdot0+0\cdot0=0$ ; colonne $C$ :
$0\cdot1+1\cdot1+1\cdot0+0\cdot0=1$ ; colonne $D$ : $0\cdot0+1\cdot1+1\cdot1+0\cdot0=2$.
2) Le terme $(A,D)$ de $M^2$ vaut $2$ : il y a $2$ chemins de longueur $2$ de $A$ vers $D$,
à savoir $A\to B\to D$ et $A\to C\to D$.
3) Le nombre total de chemins de longueur $2$ est la somme de tous les termes de $M^2$ :
$(0+0+1+2)+(1+0+0+1)+(1+0+0+0)+(0+1+1+0)=3+2+1+2=8$.

\noindent\textbf{Exercice 3.}
1) En $+\infty$ : $f(x)=\frac{x}{\mathrm e^{x}}\to0^+$ (croissances comparées) ; en
$-\infty$ : $f(x)\to-\infty$. $f'(x)=\mathrm e^{-x}-x\mathrm e^{-x}=(1-x)\mathrm e^{-x}$ :
positif sur $]-\infty,1[$, négatif sur $]1,+\infty[$. $f$ croît puis décroît, maximum
$f(1)=\mathrm e^{-1}$.
2) Sur $[0,1]$, $f\ge0$, donc $\mathcal A=\int_0^1 x\mathrm e^{-x}\dd x$. IPP ($u=x$,
$v'=\mathrm e^{-x}$, $v=-\mathrm e^{-x}$) :
$\mathcal A=[-x\mathrm e^{-x}]_0^1+\int_0^1\mathrm e^{-x}\dd x=-\mathrm e^{-1}+[-\mathrm e^{-x}]_0^1=-\mathrm e^{-1}-\mathrm e^{-1}+1=1-2\mathrm e^{-1}$.
Numériquement $\mathcal A\approx1-0{,}7358=0{,}264\ \si{\centi\metre\squared}$.
3) $V=\pi\int_0^1(\mathrm e^{-x})^2\dd x=\pi\int_0^1\mathrm e^{-2x}\dd x=\pi\Big[-\tfrac12\mathrm e^{-2x}\Big]_0^1
=\pi\Big(\tfrac12-\tfrac12\mathrm e^{-2}\Big)=\dfrac{\pi}{2}(1-\mathrm e^{-2})\approx1{,}358$ u.v.

\noindent\textbf{Exercice 4.}
1) $\overrightarrow{AB}=(4\,;\,2)$, donc $\overrightarrow{AB}\cdot\overrightarrow{AB}=16+4=20$
et $AB=\sqrt{20}=2\sqrt5$.
2) $\overrightarrow{MA}\cdot\overrightarrow{MB}=0$ signifie que $M$ voit $[AB]$ sous un
angle droit (ou $M\in\{A,B\}$) : $(\Gamma_1)$ est le \textbf{cercle de diamètre $[AB]$},
de centre $I(3\,;\,3)$ et de rayon $\frac{AB}{2}=\sqrt5$.
3) Avec $MA^2+MB^2=2MI^2+\frac12AB^2$ et $\frac12AB^2=\frac12\times20=10$ :
$MA^2+MB^2=40\iff2MI^2+10=40\iff MI^2=15\iff MI=\sqrt{15}$. Donc $(\Gamma_2)$ est le
\textbf{cercle de centre $I(3\,;\,3)$ et de rayon $\sqrt{15}$}.

\smallskip
\noindent\textbf{Partie B — Situation.}
\emph{Tâche 1.} Recette $R(x)=10x$ (milliers de F). Bénéfice
$B(x)=R(x)-C(x)=10x-(x^2-12x+80)=-x^2+22x-80$. $B'(x)=-2x+22=0$ en $x=11$. Le bénéfice
est maximal pour $x=11$ titres : $B(11)=-121+242-80=41$, soit \textbf{41 milliers de F CFA}.

\emph{Tâche 2.} Abonnés au mois $n$ : $a_n=2000\times1{,}05^{\,n-1}$. On résout $a_n>3000$,
soit $1{,}05^{\,n-1}>1{,}5$, d'où $n-1>\frac{\ln1{,}5}{\ln1{,}05}\approx8{,}31$ : à partir
du \textbf{10\textsuperscript{e} mois} ($a_{10}=2000\times1{,}05^9\approx3102$ abonnés).

\emph{Tâche 3.} $P(\text{au moins une erreur})=1-(0{,}92)^3=1-0{,}778688=\mathbf{0{,}221312}$,
soit environ \SI{22}{\percent}.
